\section{The parameterized semiprime denominator system}
\label{sec:denominator-system}

Fix throughout a reduced target
\[
 t=\frac ab\in(0,1)
\]
with squarefree $b$, a fixed anchor ratio $\eta\in(0,1)$, a finite forbidden set $\mathcal T$, and a fixed exponent $\gamma>1$. Reducedness and $0<t<1$ imply $b>1$. Thus the set $S_b$ of prime divisors of $b$ is nonempty and
\begin{equation}
\label{eq:tau-positive}
 \tau(b):=\sum_{r\in S_b}\frac1{r^2}>0.
\end{equation}
Put
\begin{equation}
\label{eq:lambda-gamma}
 \lambda_\gamma=\frac{(\log\gamma)^2}{2}.
\end{equation}
We assume
\begin{equation}
\label{eq:admissible-region}
 t<\lambda_\gamma<1,
\end{equation}
equivalently $\exp(\sqrt{2t})<\gamma<\exp(\sqrt2)$. All asymptotic statements are for $X\to\infty$ with $t,\gamma,\eta,b$ and $\mathcal T$ fixed. The estimates are uniform when $\eta$ ranges in a compact subset of $(0,1)$.

\subsection{Prime blocks and reciprocal masses}

Put
\begin{equation}
\label{eq:prime-blocks}
 Z=X^\gamma,
 \qquad P=\{p\text{ prime}:X\leq p<Z\},
 \qquad B=\{q\text{ prime}:\eta Z\leq q<Z\}.
\end{equation}
We take $X$ so large that every prime in $P$ exceeds every prime divisor of $b$, and every denominator below exceeds every element of $\mathcal T$.

We use the prime number theorem and its standard partial-summation consequences~\cite[Chapter~1]{montgomery-vaughan-2007}.

\begin{lemma}[Prime supply in fixed-ratio intervals]
\label{lem:fixed-ratio-prime-supply}
For fixed $0<u<v$, as $y\to\infty$,
\[
 \#\{p\text{ prime}:uy\leq p<vy\}
 =\bigl(v-u+o_{u,v}(1)\bigr)\frac{y}{\log y}.
\]
The error is uniform on compact subsets of $0<u<v$.
\end{lemma}
\begin{proof}
Apply the prime number theorem at $uy$ and $vy$.
\end{proof}

Define
\[
 S_X=\sum_{p\in P}\frac1p,
 \quad U_X=\sum_{p\in P}\frac1{p^2},
 \quad V_B=\sum_{q\in B}\frac1q,
 \quad W_B=\sum_{q\in B}\frac1{q^2}.
\]

\begin{lemma}[Prime reciprocal estimates]
\label{lem:reciprocal-prime-estimates}
As $X\to\infty$,
\begin{align}
 S_X&=\log\gamma+o(1), & U_X&\sim\frac1{X\log X},\label{eq:prime-moments}\\
 |P|&\sim\frac{Z}{\log Z}, & |B|&=\bigl(1-\eta+o_\eta(1)\bigr)\frac{Z}{\log Z},\label{eq:B-cardinality}\\
 V_B&=\frac{\log(1/\eta)+o_\eta(1)}{\log Z}, &
 W_B&=\frac{\eta^{-1}-1+o_\eta(1)}{Z\log Z}.
 \label{eq:B-power-sums}
\end{align}
\end{lemma}
\begin{proof}
Partial summation gives
\[
 \sum_{X\leq p<X^\gamma}\frac1p
 =\int_X^{X^\gamma}\frac{du}{u\log u}+o(1)=\log\gamma+o(1)
\]
and
\[
 \sum_{X\leq p<X^\gamma}\frac1{p^2}
 =\int_X^{X^\gamma}\frac{du}{u^2\log u}(1+o(1))\sim\frac1{X\log X}.
\]
The cardinalities follow from \cref{lem:fixed-ratio-prime-supply}. The same calculation on $[\eta Z,Z)$ gives the two displayed anchor moments.
\end{proof}

\subsection{The denominator family}

Define
\[
 \mathcal E_{\rm pair}=\{pq:p,q\in P,\ p<q\},
 \qquad
 \mathcal E_b=\{rq:r\in S_b,\ q\in B\},
\]
and
\begin{equation}
\label{eq:E-and-L}
 \mathcal E=\mathcal E_{\rm pair}\sqcup\mathcal E_b,
 \qquad L=b\prod_{p\in P}p.
\end{equation}
Put $\Lambda=\sum_{e\in\mathcal E}1/e$.

\begin{proposition}[Arithmetic capacity]
\label{prop:arithmetic-capacity}
For all sufficiently large $X$, the elements of $\mathcal E$ are pairwise distinct squarefree semiprimes, they avoid $\mathcal T$, and they divide $L$. Moreover,
\begin{equation}
\label{eq:load-limit}
 \Lambda\longrightarrow\lambda_\gamma,
\end{equation}
and
\begin{equation}
\label{eq:complete-square-load}
 \sum_{e\in\mathcal E}\frac1{e^2}
 =\frac{1+o(1)}{2X^2\log^2X}
 +\frac{(\eta^{-1}-1)\tau(b)+o_{b,\eta}(1)}{Z\log Z}.
\end{equation}
In particular, $t<\Lambda<1$ for all sufficiently large $X$.
\end{proposition}
\begin{proof}
Unique factorization gives the elementary assertions. The complete-pair identity yields
\[
 \sum_{p<q\in P}\frac1{pq}=\frac12(S_X^2-U_X)\longrightarrow\frac{(\log\gamma)^2}{2}.
\]
The target-row load is
\[
 \left(\sum_{r\in S_b}\frac1r\right)V_B=o_{b,\eta}(1),
\]
proving \cref{eq:load-limit}. Applying the same identity to $u_p=p^{-2}$ gives
\[
 \sum_{p<q\in P}\frac1{p^2q^2}
 =\frac12\left(U_X^2-\sum_{p\in P}\frac1{p^4}\right)
 \sim\frac1{2X^2\log^2X}.
\]
The target rows contribute $\tau(b)W_B$, which gives \cref{eq:complete-square-load}.
\end{proof}

\subsection{Exact centring and the actual variance}

Set
\begin{equation}
\label{eq:theta-definition}
 \theta=\frac{t}{\Lambda}.
\end{equation}
Then $\theta$ lies in a fixed compact subinterval $I_{t,\gamma,\eta}\subset(0,1)$ and
\begin{equation}
\label{eq:theta-limit}
 \theta\longrightarrow\alpha_{t,\gamma}:=\frac{2t}{(\log\gamma)^2}.
\end{equation}
Moreover,
\begin{equation}
\label{eq:exact-centering-load}
 \theta\sum_{e\in\mathcal E}\frac1e=t,
 \qquad \Lambda<1.
\end{equation}
Define
\begin{equation}
\label{eq:sigma-E}
 \sigma_{\mathcal E}^2=\theta(1-\theta)\sum_{e\in\mathcal E}\frac1{e^2}.
\end{equation}

\begin{proposition}[Total actual-family variance]
\label{prop:variance-regimes}
With $\alpha=\alpha_{t,\gamma}$,
\begin{equation}
\label{eq:total-variance}
 \sigma_{\mathcal E}^2\sim\alpha(1-\alpha)
 \left(\frac1{2X^2\log^2X}+\frac{(\eta^{-1}-1)\tau(b)}{Z\log Z}\right).
\end{equation}
\end{proposition}
\begin{proof}
Combine \cref{eq:theta-limit,eq:complete-square-load,eq:sigma-E}.
\end{proof}

The prime-pair term dominates when $\gamma>2$; the target-row term dominates when $1<\gamma\leq2$. The proof below always uses the sum in \cref{eq:total-variance}, so no separate argument is required at the boundary.

The finite Fourier numerator is
\begin{equation}
\label{eq:Fh-definition}
 F(h)=\exp\left(-\frac{2\pi i a h}{b}\right)
 \prod_{e\in\mathcal E}\left((1-\theta)+\theta\exp\left(\frac{2\pi i h}{e}\right)\right).
\end{equation}
Compactness of $I_{t,\gamma,\eta}$ gives $\kappa_{\rm mod}>0$ such that
\begin{equation}
\label{eq:bernoulli-modulus}
 |(1-\theta)+\theta\exp(2\pi i u)|
 \leq\exp(-\kappa_{\rm mod}\dist{u}^2).
\end{equation}

\subsection{Exact factor partition}

Chinese-remainder coordinates split the frequency into those indexed by $B$, by $P\setminus B$, and by $S_b$. The denominator factors are assigned exactly once:
\begin{center}
\begin{tabular}{ll}
\toprule
Denominator class & use \\
\midrule
$q q'$ with $q,q'\in B$ & internal anchor factor \\
$r q$ with $r\in P\setminus B$, $q\in B$ & lower-prime row $r$ \\
$r q$ with $r\in S_b$, $q\in B$ & target-coordinate row $r$ \\
$p p'$ with $p,p'\in P\setminus B$ & retained prime-pair factor \\
\bottomrule
\end{tabular}
\end{center}
The target character and every phase not absorbed into a row remain in the complex residual factor.
